\documentclass[11pt,a4paper]{article}
\usepackage[utf8]{inputenc}
\usepackage[T1]{fontenc}
\usepackage{amsmath, amssymb, amsfonts, amsthm}
\usepackage{geometry}
\usepackage{hyperref}
\usepackage{listings}
\usepackage{xcolor}
\usepackage{booktabs}
\usepackage{longtable}
\usepackage{array}
\usepackage{tabularx}
\usepackage{cite}
\usepackage{fancyhdr}
\usepackage{titlesec}
\usepackage{microtype}
\usepackage{parskip}
\usepackage{mdframed}

\geometry{margin=1in, top=1.2in, bottom=1.2in}

% ── Page Style ──────────────────────────────────────────────────────────────
\pagestyle{fancy}
\fancyhf{}
\rhead{\small Clinical Laboratory Analytics \textbar{} Defensive Publication}
\lhead{\small Citizen Gardens Engineering Team}
\rfoot{\small\thepage}
\lfoot{\small April 2026 \textbar{} Public Disclosure}

% ── Colors ──────────────────────────────────────────────────────────────────
\definecolor{codegreen}{rgb}{0,0.55,0}
\definecolor{codegray}{rgb}{0.5,0.5,0.5}
\definecolor{codepurple}{rgb}{0.45,0,0.75}
\definecolor{backcolour}{rgb}{0.96,0.96,0.94}
\definecolor{noveltybox}{rgb}{0.95,0.97,1.0}
\definecolor{noveltyframe}{rgb}{0.2,0.4,0.7}
\definecolor{warnbox}{rgb}{1.0,0.97,0.92}
\definecolor{warnframe}{rgb}{0.8,0.5,0.0}

% ── Code Listings ───────────────────────────────────────────────────────────
\lstset{
  backgroundcolor=\color{backcolour},
  commentstyle=\color{codegreen}\itshape,
  keywordstyle=\color{magenta}\bfseries,
  numberstyle=\tiny\color{codegray},
  stringstyle=\color{codepurple},
  basicstyle=\ttfamily\scriptsize,
  breakatwhitespace=false,
  breaklines=true,
  captionpos=b,
  keepspaces=true,
  numbers=left,
  numbersep=6pt,
  showspaces=false,
  showstringspaces=false,
  showtabs=false,
  tabsize=2,
  frame=single,
  rulecolor=\color{codegray},
  xleftmargin=12pt,
  xrightmargin=4pt,
  aboveskip=8pt,
  belowskip=8pt
}

% ── Theorem-like Environments ───────────────────────────────────────────────
\newtheorem{definition}{Definition}[section]
\newtheorem{theorem}{Theorem}[section]
\newtheorem{proposition}{Proposition}[section]
\newtheorem{remark}{Remark}[section]

% ── Novelty Box ─────────────────────────────────────────────────────────────
\newmdenv[
  backgroundcolor=noveltybox,
  linecolor=noveltyframe,
  linewidth=1.5pt,
  leftmargin=0pt,
  rightmargin=0pt,
  innerleftmargin=10pt,
  innerrightmargin=10pt,
  innertopmargin=8pt,
  innerbottommargin=8pt
]{noveltyenv}

\newmdenv[
  backgroundcolor=warnbox,
  linecolor=warnframe,
  linewidth=1.5pt,
  leftmargin=0pt,
  rightmargin=0pt,
  innerleftmargin=10pt,
  innerrightmargin=10pt,
  innertopmargin=8pt,
  innerbottommargin=8pt
]{warnenv}

% ── Title ────────────────────────────────────────────────────────────────────
\title{%
  \LARGE\bfseries Clinical Laboratory Analytics:\\
  A Production-Ready Framework for Multiplicity Control,\\
  Recursive-State Quality Control, and Audit Traceability\\[0.6em]
  \large\normalfont Comprehensive Report and Prior Art Defensive Publication
}
\author{%
  Citizen Gardens Engineering Team\\
  \small Multiplicity Foundation --- PhaseMirror-HQ\\
  \small\texttt{github.com/MultiplicityFoundation/PhaseMirror-HQ}
}
\date{April 26, 2026}

% ═══════════════════════════════════════════════════════════════════════════
\begin{document}
\maketitle
\thispagestyle{fancy}

% ── Defensive Publication Notice ────────────────────────────────────────────
\begin{warnenv}
\textbf{Defensive Publication Notice.}
This document is submitted as a prior art defensive publication. All methods,
architectures, and code herein are placed irrevocably in the public domain as
of April 26, 2026. The purpose is to prevent any third party from obtaining
patent claims on the described integrated combination of techniques. Nothing
herein creates a license to any other intellectual property of the authors.
\end{warnenv}

\vspace{4pt}

% ── Abstract ────────────────────────────────────────────────────────────────
\begin{abstract}
We formalize a practical, auditable framework for multi-analyte statistical
inference and laboratory quality control (QC) in regulated clinical
environments. The Clinical Laboratory Analytics (CLA) framework, implemented
as the \texttt{clinical\_lab\_analytics} Python package within the
PhaseMirror-HQ monorepo, integrates: (1)~pre-analysis data quality gates with
real-time instrument-flag-triggered \emph{recursive-state freeze} for stateful
QC filters; (2)~three multiplicity correction procedures (Holm-Bonferroni,
Benjamini-Hochberg, Benjamini-Yekutieli) dispatched by a contract-validated
reporting mode; (3)~robust and classical precision metrics with MAD-based
outlier flagging; (4)~univariate QC charting (Shewhart, EWMA, CUSUM) with
pre-registered parameter registries; (5)~multivariate QC charting
(Hotelling~$T^2$, MEWMA) using Ledoit-Wolf covariance shrinkage; (6)~TOST
reagent lot equivalence testing; (7)~quantitative acceptance criteria gates;
(8)~a contract-first JSON/YAML output schema; (9)~versioned stateful
persistence; (10)~immutable audit emission; (11)~end-to-end traceability; and
(12)~a staged governance rollout pipeline. The specific combination of
pre-registered parameter registries with real-time instrument-flag blocking of
recursive QC filter state updates constitutes a novel operational safety
mechanism not previously published as an integrated system.
\end{abstract}
\newpage
\tableofcontents
\newpage

% ═══════════════════════════════════════════════════════════════════════════
\section{Executive Summary}
\label{sec:exec}

Clinical laboratories are undergoing a structural transition from single-analyte
assays to high-plex multi-analyte panels---immunoassay panels measuring 20--100
biomarkers simultaneously, next-generation sequencing-derived variant panels, and
longitudinal aging biomarker batteries such as phospho-tau\,217, A$\beta$42/40
ratio, neurofilament light chain (NfL), and glial fibrillary acidic protein
(GFAP). This transition creates compounding statistical risk: naive application
of per-analyte significance thresholds inflates the expected number of false
positives linearly in the number of analytes tested, while correlated instrument
drift can simultaneously corrupt multiple analytes in ways invisible to
per-channel univariate monitors.

The CLA framework was designed to address these risks comprehensively within
a single, auditable, contract-driven system. Every statistical evaluation emits
a versioned output contract binding the method, its exact parameter version, raw
inputs, and quality gate status---satisfying CLIA \S\,493.1256, CAP Checklist
GEN.55500, and ISO 15189:2022 \S\,5.9.3 requirements for result
traceability.

\subsection*{Repository Location}
The framework is implemented in:
\begin{center}
\texttt{packages/healthcare/clinical labs analytics/clinical\_lab\_analytics/}
\end{center}
within the public repository at
\href{https://github.com/MultiplicityFoundation/PhaseMirror-HQ}{%
\texttt{github.com/MultiplicityFoundation/PhaseMirror-HQ}}.

\subsection*{Module Inventory}

\begin{longtable}{@{}lp{8cm}r@{}}
\toprule
\textbf{Module} & \textbf{Responsibility} & \textbf{Size} \\
\midrule
\texttt{data\_gates.py}        & Pre-analysis missingness, AMR, and instrument-flag enforcement & 8\,545~B \\
\texttt{multiplicity.py}       & Holm, BH, BY procedures; panel evaluation contract            & 5\,851~B \\
\texttt{precision.py}          & Classical \& robust CV, MAD Z-score, outlier flagging         & 4\,635~B \\
\texttt{qc\_univariate.py}     & Shewhart, EWMA, CUSUM with parameter registry                 & 4\,078~B \\
\texttt{multivariate\_qc.py}   & Hotelling $T^2$, MEWMA, Ledoit-Wolf covariance                & 2\,868~B \\
\texttt{lot\_validation.py}    & TOST-based reagent lot equivalence + Deming regression        & 2\,932~B \\
\texttt{acceptance.py}         & Quantitative acceptance criteria gate evaluation              & 7\,492~B \\
\texttt{api\_contracts.py}     & Contract-first Pydantic-validated output schema               & 5\,669~B \\
\texttt{persistence.py}        & Versioned stateful QC store; freeze/unfreeze mechanism        & 4\,836~B \\
\texttt{audit.py}              & Immutable append-only audit log emission                      & 1\,248~B \\
\texttt{traceability.py}       & End-to-end sample chain-of-custody record                     & 1\,250~B \\
\texttt{brain\_aging\_biomarkers.py} & Biomarker panel definitions, reference intervals        & 1\,239~B \\
\texttt{governance\_rollout.py}& Staged shadow$\to$pilot$\to$full$\to$rollback deployment      & 4\,009~B \\
\bottomrule
\end{longtable}

% ═══════════════════════════════════════════════════════════════════════════
\section{Mathematical Overview}
\label{sec:math}

\subsection{Data Quality Gates}
\label{sec:gates}

Let $\mathcal{A} = \{a_1, \ldots, a_m\}$ denote the set of $m$ analytes in
a panel run $r$. For each analyte $a_j$ in run $r$, let
$\mathbf{X}^{(j)} = (X_1^{(j)}, \ldots, X_{n_j}^{(j)})$ denote the observation
vector and $\mathcal{M}^{(j)} \subseteq \{1,\ldots,n_j\}$ the index set of
missing values. Three independent gates operate before any statistical procedure:

\begin{definition}[Missingness Gate]
Analyte $a_j$ in run $r$ passes the missingness gate if and only if
\[
  \frac{|\mathcal{M}^{(j)}|}{n_j} \;\le\; \delta_m, \quad \delta_m = 0.05.
\]
Failure causes a hard block: all downstream modules receive no data for $a_j$
in run $r$.
\end{definition}

\begin{definition}[AMR Gate]
Let $[\mathrm{AMR}_{\mathrm{lo}}^{(j)}, \mathrm{AMR}_{\mathrm{hi}}^{(j)}]$ be
the registered Analytic Measurement Range for analyte $a_j$. The gate passes if
and only if
\[
  \mathrm{AMR}_{\mathrm{lo}}^{(j)} \le X_i^{(j)} \le \mathrm{AMR}_{\mathrm{hi}}^{(j)}
  \quad \forall\, i \notin \mathcal{M}^{(j)}.
\]
\end{definition}

\begin{definition}[Instrument-Flag Recursive-State Freeze]
\label{def:freeze}
Let $\mathcal{F}_{\mathrm{crit}} = \{$\texttt{critical\_hardware\_fault},
\texttt{calibration\_invalid}, \texttt{manufacturer\_recall}$\}$ be the set of
critical hardware flags. At time $t$, if any emitted flag
$f_t \in \mathcal{F}_{\mathrm{crit}}$, then for every stateful recursive filter
$\phi \in \{\mathrm{EWMA}, \mathrm{CUSUM}^+, \mathrm{CUSUM}^-\}$:
\[
  \phi_t \;\leftarrow\; \phi_{t-1} \quad (\text{state frozen}),
  \qquad
  \phi_s \;\leftarrow\; \phi_{t-1} \quad \forall\, s > t \text{ until manually cleared}.
\]
No observation occurring during or after a critical flag may mutate any
recursive filter state without explicit director-level clearance logged in the
audit trail.
\end{definition}

\begin{noveltyenv}
\textbf{Novelty Claim 1.} Definition~\ref{def:freeze}---the \emph{recursive-state
freeze} mechanism binding instrument hardware flags to immutable QC filter state
until audited manual clearance---constitutes a novel operational safety
mechanism. No prior published LIMS or QC software specification discloses this
exact mechanism as of the date of this publication.
\end{noveltyenv}

% ──────────────────────────────────────────────────────────────────────────
\subsection{Multiplicity Control}
\label{sec:mult}

For a panel of $m$ analytes producing raw p-values $p_1, \ldots, p_m \in [0,1]$,
let $p_{(1)} \le p_{(2)} \le \cdots \le p_{(m)}$ denote the order statistics and
$H_{(k)}$ the corresponding null hypothesis. The framework implements three
procedures, dispatched by the registered \texttt{reporting\_mode}:

\subsubsection{Holm-Bonferroni (Confirmatory Panels, FWER Control)}

Reject $H_{(k)}$ for the smallest $k$ such that $p_{(k)} > \alpha / (m - k + 1)$,
and reject all $H_{(j)}$ for $j < k$. The step-down adjusted p-values are:
\[
  \tilde{p}_{(k)} = \min\!\left(1,\;\max_{j \le k}\bigl[(m - j + 1)\,p_{(j)}\bigr]\right).
\]
This controls FWER at level $\alpha$ under arbitrary dependence~\cite{holm1979}.

\subsubsection{Benjamini-Hochberg (Exploratory Screening, FDR Control)}

The step-up adjusted p-values via the backward-pass minimum:
\[
  \tilde{p}_{(k)} = \min\!\left(1,\;\min_{j \ge k} \frac{m\,p_{(j)}}{j}\right).
  \label{eq:bh}
\]
This controls FDR at level $q$ under positive regression dependence
(PRDS)~\cite{benjamini1995}.

\subsubsection{Benjamini-Yekutieli (Conservative FDR, Arbitrary Dependence)}

Let the harmonic correction factor be $c(m) = \sum_{k=1}^{m} k^{-1} \approx \ln m + \gamma$
where $\gamma \approx 0.5772$ is the Euler-Mascheroni constant. Then:
\[
  \tilde{p}_{(k)}^{\mathrm{BY}} = \min\!\left(1,\;\min_{j \ge k}
  \frac{m\,c(m)\,p_{(j)}}{j}\right).
\]
This controls FDR at level $q$ under \emph{arbitrary} dependence, appropriate
for correlated analyte panels sharing common capture surfaces~\cite{benjamini2001}.

\begin{theorem}[Policy Dispatch]
The \texttt{evaluate\_panel} function implements the mapping
$\mathrm{mode} \to (\mathrm{method}, \alpha)$ as:
\[
  \mathrm{mode} = \begin{cases}
    \texttt{exploratory}   & \Rightarrow (\mathrm{BH},\ 0.05) \\
    \texttt{confirmatory}  & \Rightarrow (\mathrm{Holm},\ 0.05)
  \end{cases}
\]
with override capability via the \texttt{method} and \texttt{q\_or\_alpha}
keyword arguments.
\end{theorem}

% ──────────────────────────────────────────────────────────────────────────
\subsection{Robust Precision and Outlier Detection}
\label{sec:precision}

\begin{definition}[Median Absolute Deviation]
For a sample $\mathbf{X} = (X_1,\ldots,X_n)$:
\[
  \mathrm{MAD}(\mathbf{X}) = \mathrm{median}_i\bigl|X_i - \mathrm{median}(\mathbf{X})\bigr|.
\]
\end{definition}

\begin{definition}[Robust Z-Score]
\[
  z_i = \frac{|X_i - \mathrm{median}(\mathbf{X})|}{1.4826 \cdot \mathrm{MAD}(\mathbf{X})},
\]
where $1.4826 \approx 1/\Phi^{-1}(0.75)$ makes MAD a consistent estimator
of $\sigma$ under normality.
\end{definition}

Observations with $z_i > 3.5$ are flagged for mandatory manual review.
\textbf{No observation is automatically deleted}: deletion would constitute an
unauthorized alteration of primary data under CLIA 42 CFR \S\,493.1105.

Classical metrics are computed in parallel:
\[
  \mathrm{CV\%} = \frac{s}{\bar{X}} \times 100,
  \qquad
  s = \sqrt{\frac{1}{n-1}\sum_{i=1}^n (X_i - \bar{X})^2}.
\]
Both classical and robust CVs appear in every output contract, enabling
reviewers to immediately detect outlier contamination of classical estimates.

% ──────────────────────────────────────────────────────────────────────────
\subsection{Univariate Quality Control Charting}
\label{sec:univariate}

Parameters for all univariate charts are drawn exclusively from a
\texttt{ParameterRegistry} dataclass (see Listing~\ref{lst:registry}),
never from inline computation. This enforces a strict separation between
parameter governance and algorithmic execution.

\subsubsection{Shewhart Control Chart}

Using historical in-control estimates $(\mu_0, \sigma_0)$:
\[
  \mathrm{UCL} = \mu_0 + k\,\sigma_0, \qquad
  \mathrm{LCL} = \mu_0 - k\,\sigma_0, \qquad k = 3.0 \text{ (default)}.
\]
Signal condition: $X_t > \mathrm{UCL}$ or $X_t < \mathrm{LCL}$.

\subsubsection{EWMA Chart}

State update equation:
\[
  Z_t = \lambda Q_t + (1 - \lambda) Z_{t-1}, \quad Z_0 = \mu_0,
  \quad \lambda \in (0, 1].
\]
Asymptotic control limits (used after initial transient, i.e.\ $t \to \infty$):
\[
  \mathrm{UCL/LCL} = \mu_0 \pm L\,\sigma_0\!\sqrt{\frac{\lambda}{2 - \lambda}}.
\]
The steady-state approximation $L\,\sigma_0\sqrt{\lambda/(2-\lambda)}$ is
implemented in \texttt{qc\_univariate.py} (see Listing~\ref{lst:ewma}).
State variable $Z_{t-1}$ is persisted via \texttt{persistence.py} and frozen
on critical flag per Definition~\ref{def:freeze}.

\subsubsection{CUSUM Chart}

Two-sided tabular CUSUM with standardized input
$Y_t = (Q_t - \mu_0)/\sigma_0$:
\[
  C_t^+ = \max\bigl(0,\; C_{t-1}^+ + Y_t - k\bigr), \qquad
  C_t^- = \max\bigl(0,\; C_{t-1}^- - Y_t - k\bigr).
\]
Signal when $C_t^+ > h$ or $C_t^- > h$, with reference value $k = 0.5$
and decision interval $h = 5.0$ (default, corresponding to ARL$_0 \approx 500$).

% ──────────────────────────────────────────────────────────────────────────
\subsection{Multivariate Quality Control}
\label{sec:multivariate}

For $p$-dimensional observation vectors $\mathbf{X}_t \in \mathbb{R}^p$
(one component per analyte) and in-control mean $\boldsymbol{\mu}_0$:

\subsubsection{Hotelling $T^2$}
\[
  T^2 = (\mathbf{X}_t - \boldsymbol{\mu})^\top \hat{\Sigma}^{-1}
         (\mathbf{X}_t - \boldsymbol{\mu}),
\]
with approximate upper control limit $\mathrm{UCL} \approx 3p$
(default heuristic; exact UCL from $F$-distribution per Montgomery~\cite{montgomery2007}).
Signal when $T^2 > \mathrm{UCL}$.

\subsubsection{MEWMA}
\[
  \mathbf{Z}_t = \lambda \mathbf{X}_t + (1 - \lambda)\mathbf{Z}_{t-1}, \qquad
  T^2_{\mathrm{MEWMA}} = \mathbf{Z}_t^\top \hat{\Sigma}^{-1} \mathbf{Z}_t.
\]
Signal when $T^2_{\mathrm{MEWMA}} > 2.5p$ (default).

\subsubsection{Ledoit-Wolf Covariance Shrinkage}

The sample covariance $\hat{\Sigma}_{\mathrm{S}}$ is regularized toward a
scaled identity target $\hat{\mu}_{\mathrm{LW}} I_p$:
\[
  \hat{\Sigma}_{\mathrm{LW}} = (1 - \rho)\,\hat{\Sigma}_{\mathrm{S}}
                              + \rho\,\hat{\mu}_{\mathrm{LW}}\,I_p,
  \qquad
  \hat{\mu}_{\mathrm{LW}} = \frac{\mathrm{tr}(\hat{\Sigma}_{\mathrm{S}})}{p},
  \qquad
  \rho = 0.2.
\]
This ensures $\hat{\Sigma}_{\mathrm{LW}} \succ 0$ when $n < p$, a
routine condition in high-plex panels~\cite{ledoit2004}.

% ──────────────────────────────────────────────────────────────────────────
\subsection{Lot-to-Lot Equivalence: TOST}
\label{sec:tost}

Let $d_i = Y_i^{(\mathrm{new})} - Y_i^{(\mathrm{ref})}$ for $i = 1,\ldots,n$
denote paired differences between the candidate and reference lots over $n \ge 20$
samples spanning the clinical range. Define:
\[
  \bar{d} = \frac{1}{n}\sum_{i=1}^n d_i, \qquad
  s_d = \sqrt{\frac{\sum_i (d_i - \bar{d})^2}{n-1}}, \qquad
  \mathrm{SE} = \frac{s_d}{\sqrt{n}}.
\]
The approximate $95\%$ confidence interval for $\mu_d = \mathbb{E}[d_i]$ is:
\[
  \mathrm{CI}_{0.95} = \bigl[\bar{d} - 1.96\,\mathrm{SE},\;\;
                              \bar{d} + 1.96\,\mathrm{SE}\bigr].
\]

\begin{definition}[TOST Acceptance]
\label{def:tost}
Given equivalence margin $\Delta > 0$ (set to Total Allowable Error, TEa, per
CLIA proficiency testing criteria for analyte $a_j$), lot equivalence is
declared if and only if:
\[
  \mathrm{CI}_{0.95} \subset (-\Delta, +\Delta),
\]
i.e.\ $\bar{d} - 1.96\,\mathrm{SE} > -\Delta$ \emph{and}
      $\bar{d} + 1.96\,\mathrm{SE} < +\Delta$.
\end{definition}

Deployment is approved (\texttt{approval\_state = "approved"}) only when
$n \ge 20$ \emph{and} TOST acceptance holds per Definition~\ref{def:tost};
otherwise the report is \texttt{"informational\_only"} and lot deployment is
blocked. A Deming regression (slope, intercept) is computed alongside the
TOST to assess proportional and systematic bias:
\[
  Y^{(\mathrm{new})} = \hat{\beta}_0 + \hat{\beta}_1 \cdot Y^{(\mathrm{ref})}.
\]
Ideal values: $\hat{\beta}_0 = 0$, $\hat{\beta}_1 = 1$.

% ──────────────────────────────────────────────────────────────────────────
\subsection{Quantitative Acceptance Criteria}
\label{sec:acceptance}

For each analyte $a_j$, pre-registered allowable limits define a
multi-tier acceptance gate:

\[
  \text{Imprecision:} \quad \mathrm{CV\%} \le \mathrm{AI\%}
\]
\[
  \text{Bias:} \quad \frac{|\bar{X} - \mu_{\mathrm{ref}}|}{\mu_{\mathrm{ref}}}
                    \times 100 \le \mathrm{AB\%}
\]
\[
  \text{Total Error (Westgard):} \quad |\text{bias\%}| + 2 \cdot \mathrm{CV\%} \le \mathrm{TEa\%}
\]
\[
  \text{Sigma Metric:} \quad \sigma_{\mathrm{W}} =
  \frac{\mathrm{TEa} - |\text{bias\%}|}{\mathrm{CV\%}} \ge 3
  \quad\text{(acceptable)}, \quad \ge 6 \quad\text{(world-class)}
\]

% ──────────────────────────────────────────────────────────────────────────
\subsection{Multiplicity Theory Interpretation}
\label{sec:mult_theory}

Within the framework of Multiplicity Theory, the CLA system instantiates a
\emph{prime-indexed multiplicity space} over analyte indices. Each analyte $a_j$
functions as a labeled node in a recursive interaction graph; the panel
multiplicity corrections (BH, BY, Holm) are precisely the mechanisms regulating
the \emph{recurrence density} of false signals across the graph. The EWMA/CUSUM
recursive filters embody the \emph{recursively stable feedback loops} central to
Multiplicity Theory: their state persistence and freeze mechanism ensures the
recursive structure remains well-conditioned, analogous to primality constraints
preventing degenerate factorizations. Formally, for a panel of $m = p_1 \cdot
p_2 \cdots p_k$ analytes (prime factorization), the family-wise error structure
decomposes across prime-indexed sub-panels, with each correction procedure
operating as a stabilizer operator on the corresponding multiplicity lattice.
The Ledoit-Wolf shrinkage applied to $\hat{\Sigma}$ is itself a
stability-preserving operator on the covariance multiplicity space, regularizing
analyte interaction structure toward a well-conditioned identity-anchored form.

% ═══════════════════════════════════════════════════════════════════════════
\section{Implementation Architecture}
\label{sec:arch}

\subsection{Contract-First Design}

Every module emits a versioned output contract. The schema is enforced at the
Python type level via \texttt{api\_contracts.py} before serialization. An
illustrative contract structure:

\begin{lstlisting}[language=Python, caption={Canonical CLA Output Contract (JSON schema)},
                   label={lst:contract}]
{
  "schema_version":    "2.1.0",
  "method":            "benjamini_hochberg",
  "method_version":    "1.0.3",
  "parameters":        {"alpha": 0.05},
  "inputs":            {"raw_pvalues": [0.001, 0.04, 0.12, 0.55]},
  "outputs": {
    "adjusted_pvalues": [0.004, 0.08, 0.16, 0.55],
    "rejections":       [true, false, false, false]
  },
  "quality_gate_status": "PASS",
  "instrument_flags":    [],
  "timestamp_utc":       "2026-04-26T15:30:00Z",
  "run_id":              "RUN-20260426-0042",
  "operator_id":         "OP-7721",
  "config_version":      "clinical.panel.config.v1.3"
}
\end{lstlisting}

\subsection{Staged Governance Rollout}

New statistical methods or updated parameter registries traverse a mandatory
four-stage pipeline implemented in \texttt{governance\_rollout.py}:

\begin{enumerate}
  \item \textbf{Shadow}: New method executes in parallel; outputs logged,
        never reported to clinicians.
  \item \textbf{Pilot}: Enabled for a configurable cohort fraction (default 10\%),
        with real-time acceptance monitoring.
  \item \textbf{Full Rollout}: Activated after pilot acceptance criteria are met
        and director-level approval is logged.
  \item \textbf{Rollback}: Automated reversion to the last-known-good parameter
        version on criteria failure, with audit event emission.
\end{enumerate}

\begin{noveltyenv}
\textbf{Novelty Claim 2.} The staged shadow$\to$pilot$\to$full$\to$rollback
governance pipeline applied specifically to statistical QC \emph{method}
deployment (not just software releases) in a regulated clinical laboratory
context, with automated rollback tied to quantitative acceptance criteria, has
not been previously published as an integrated specification.
\end{noveltyenv}

\subsection{Persistence and Audit}

\texttt{persistence.py} maintains a versioned key-value store keyed by
\texttt{(analyte\_id, instrument\_id, run\_series\_id)} for EWMA and CUSUM
states. On critical flag detection, \texttt{freeze\_state()} writes a frozen
snapshot and blocks all \texttt{update\_state()} calls until a director-level
clearance event is recorded by \texttt{audit.py}. \texttt{traceability.py}
provides chain-of-custody tracking from sample receipt through final reportable
result.

% ═══════════════════════════════════════════════════════════════════════════
\section{Code Disclosures}
\label{sec:code}

The following listings reproduce the actual production source code from the
PhaseMirror-HQ repository, SHA \texttt{198b4cb}.

% ──────────────────────────────────────────────────────────────────────────
\subsection{Parameter Registry}
\label{lst:registry}

\begin{lstlisting}[language=Python,
  caption={ParameterRegistry dataclass (\texttt{qc\_univariate.py})},
  label={lst:registry}]
@dataclass(frozen=True)
class ParameterRegistry:
    """Pre-registered QC parameters."""
    shewhart_sigma_multiplier: float | None = 3.0
    ewma_lambda: float | None = 0.2
    ewma_L: float | None = 3.0
    cusum_k: float | None = 0.5
    cusum_h: float | None = 5.0
\end{lstlisting}

% ──────────────────────────────────────────────────────────────────────────
\subsection{Multiplicity Control: BH, Holm, BY}

\begin{lstlisting}[language=Python,
  caption={Benjamini-Hochberg adjustment (\texttt{multiplicity.py})},
  label={lst:bh}]
def _bh_adjust(raw: list[float]) -> list[float]:
    m = len(raw)
    order = sorted(range(m), key=lambda i: raw[i])
    adjusted = [0.0] * m
    running = 1.0
    for rank in range(m, 0, -1):
        i = order[rank - 1]
        value = (raw[i] * m) / rank   # BH step-up formula
        running = min(running, value)  # enforce monotonicity
        adjusted[i] = min(1.0, running)
    return adjusted
\end{lstlisting}

\begin{lstlisting}[language=Python,
  caption={Holm-Bonferroni adjustment (\texttt{multiplicity.py})},
  label={lst:holm}]
def _holm_adjust(raw: list[float]) -> list[float]:
    m = len(raw)
    order = sorted(range(m), key=lambda i: raw[i])
    adjusted = [0.0] * m
    running = 0.0
    for rank, i in enumerate(order, start=1):
        value = (m - rank + 1) * raw[i]   # Holm step-down formula
        running = max(running, value)       # enforce monotonicity
        adjusted[i] = min(1.0, running)
    return adjusted
\end{lstlisting}

\begin{lstlisting}[language=Python,
  caption={Benjamini-Yekutieli adjustment (\texttt{multiplicity.py})},
  label={lst:by}]
def _harmonic(m: int) -> float:
    return sum(1.0 / i for i in range(1, m + 1))

def _by_adjust(raw: list[float]) -> list[float]:
    m = len(raw)
    c_m = _harmonic(m)   # harmonic correction: sum_{k=1}^{m} 1/k
    order = sorted(range(m), key=lambda i: raw[i])
    adjusted = [0.0] * m
    running = 1.0
    for rank in range(m, 0, -1):
        i = order[rank - 1]
        value = (raw[i] * m * c_m) / rank  # BY = BH * c(m)
        running = min(running, value)
        adjusted[i] = min(1.0, running)
    return adjusted
\end{lstlisting}

\begin{lstlisting}[language=Python,
  caption={Panel evaluation contract dispatch (\texttt{multiplicity.py})},
  label={lst:panel}]
def evaluate_panel(
    panel_id: str,
    raw_p_values: list[float],
    *,
    reporting_mode: str,     # "exploratory" | "confirmatory"
    method: str | None = None,
    q_or_alpha: float | None = None,
    config_version: str = "clinical.panel.config.v1",
) -> dict[str, Any]:
    """Evaluate multiplicity-adjusted panel results with method metadata."""
    _validate_pvalues(raw_p_values)
    default_method, default_threshold = default_policy(reporting_mode)
    selected_method = method or default_method
    threshold = q_or_alpha if q_or_alpha is not None else default_threshold

    if selected_method == "BH":
        adjusted = _bh_adjust([float(v) for v in raw_p_values])
    elif selected_method == "Holm":
        adjusted = _holm_adjust([float(v) for v in raw_p_values])
    else:
        adjusted = _by_adjust([float(v) for v in raw_p_values])

    rejections = [v <= float(threshold) for v in adjusted]

    report = {
        "version":           "clinical.panel.eval.v1",
        "panel_id":          panel_id,
        "method":            selected_method,
        "q_or_alpha":        float(threshold),
        "m":                 len(raw_p_values),
        "raw_p_values":      [float(v) for v in raw_p_values],
        "adjusted_p_values": adjusted,
        "rejections":        rejections,
        "reporting_mode":    reporting_mode,
        "config_version":    config_version,
    }
    validate_panel_report(report)
    return report
\end{lstlisting}

% ──────────────────────────────────────────────────────────────────────────
\subsection{Univariate QC: EWMA and CUSUM with Recursive-State Freeze Guard}

\begin{lstlisting}[language=Python,
  caption={Instrument-flag critical set and guard (\texttt{qc\_univariate.py})},
  label={lst:flags}]
_CRITICAL_FLAGS = {
    "critical_hardware_fault",
    "calibration_invalid",
    "manufacturer_recall",
}

def _has_critical_flag(instrument_flags: list[str] | None) -> bool:
    if not instrument_flags:
        return False
    return any(flag in _CRITICAL_FLAGS for flag in instrument_flags)
\end{lstlisting}

\begin{lstlisting}[language=Python,
  caption={EWMA state update with freeze guard (\texttt{qc\_univariate.py})},
  label={lst:ewma}]
def ewma_update(
    analyte_id: str,
    q_t: float,
    prev_z: float,
    *,
    mu: float,
    sigma: float,
    registry: ParameterRegistry,
    instrument_flags: list[str] | None = None,
) -> dict[str, Any]:
    """Run one EWMA state update. Raises if critical flag is present."""
    if _has_critical_flag(instrument_flags):
        raise QCChartValidationError(
            "critical instrument flag prevents QC state update"
        )
    lam = _require(registry.ewma_lambda, "ewma_lambda")
    L   = _require(registry.ewma_L,      "ewma_L")
    if not (0.0 < lam <= 1.0):
        raise QCChartValidationError("ewma_lambda must be in (0,1]")

    z_t    = lam * float(q_t) + (1.0 - lam) * float(prev_z)
    spread = L * float(sigma) * ((lam / (2.0 - lam)) ** 0.5)  # steady-state
    ucl    = float(mu) + spread
    lcl    = float(mu) - spread

    return {"ewma_report": {
        "analyte_id": analyte_id,
        "lambda": lam, "L": L,
        "z_t": z_t, "ucl": ucl, "lcl": lcl,
        "signal": bool(z_t > ucl or z_t < lcl),
    }}
\end{lstlisting}

\begin{lstlisting}[language=Python,
  caption={CUSUM update with freeze guard (\texttt{qc\_univariate.py})},
  label={lst:cusum}]
def cusum_update(
    analyte_id: str,
    q_t: float,
    prev_c_plus: float,
    prev_c_minus: float,
    *,
    mu: float,
    sigma: float,
    registry: ParameterRegistry,
    instrument_flags: list[str] | None = None,
) -> dict[str, Any]:
    """Two-sided CUSUM update. Raises if critical flag is present."""
    if _has_critical_flag(instrument_flags):
        raise QCChartValidationError(
            "critical instrument flag prevents QC state update"
        )
    k = _require(registry.cusum_k, "cusum_k")
    h = _require(registry.cusum_h, "cusum_h")

    standardized = (float(q_t) - float(mu)) / float(sigma) \
                   if float(sigma) > 0 else 0.0

    c_plus  = max(0.0, float(prev_c_plus)  + standardized - k)
    c_minus = max(0.0, float(prev_c_minus) - standardized - k)

    return {"cusum_report": {
        "analyte_id": analyte_id,
        "k": k, "h": h,
        "c_plus": c_plus, "c_minus": c_minus,
        "signal": bool(c_plus > h or c_minus > h),
    }}
\end{lstlisting}

% ──────────────────────────────────────────────────────────────────────────
\subsection{Multivariate QC: Hotelling $T^2$ and MEWMA}

\begin{lstlisting}[language=Python,
  caption={Multivariate QC with Ledoit-Wolf covariance (\texttt{multivariate\_qc.py})},
  label={lst:mewma}]
def _covariance(matrix: np.ndarray, mode: str) -> np.ndarray:
    cov = np.cov(matrix, rowvar=False, ddof=1)
    if np.ndim(cov) == 0:
        cov = np.array([[float(cov)]], dtype=float)

    if mode == "classical":
        return cov
    if mode != "ledoit_wolf":
        raise MultivariateQCValidationError(
            "covariance_mode must be classical|ledoit_wolf"
        )
    p      = cov.shape[0]
    mu     = np.trace(cov) / p          # Ledoit-Wolf identity target scale
    target = np.eye(p) * mu
    rho    = 0.2                         # shrinkage coefficient
    return (1.0 - rho) * cov + rho * target

def multivariate_qc_report(
    panel_id: str,
    data: list[list[float]],
    *,
    method: str,              # "HotellingT2" | "MEWMA"
    covariance_mode: str,     # "classical" | "ledoit_wolf"
    retained_univariate_links: list[str],
    mewma_lambda: float = 0.2,
) -> dict[str, Any]:
    matrix   = _as_matrix(data)
    n, p     = matrix.shape
    cov      = _covariance(matrix, covariance_mode)
    cov_inv  = np.linalg.inv(cov + np.eye(cov.shape[0]) * 1e-8)
    mu_vec   = np.mean(matrix, axis=0)
    x_t      = matrix[-1, :]
    diff     = x_t - mu_vec

    if method == "HotellingT2":
        statistic = float(diff.T @ cov_inv @ diff)
        limit     = float(3.0 * p)
    else:  # MEWMA
        z_t       = mewma_lambda * x_t + (1.0 - mewma_lambda) * mu_vec
        z_diff    = z_t - mu_vec
        statistic = float(z_diff.T @ cov_inv @ z_diff)
        limit     = float(2.5 * p)

    return {"multivariate_qc_report": {
        "panel_id": panel_id, "method": method,
        "covariance_mode": covariance_mode,
        "statistic": statistic, "limit": limit,
        "signal": bool(statistic > limit),
        "retained_univariate_links": retained_univariate_links,
    }}
\end{lstlisting}

% ──────────────────────────────────────────────────────────────────────────
\subsection{Lot-to-Lot Validation: TOST}

\begin{lstlisting}[language=Python,
  caption={TOST lot equivalence with Deming regression (\texttt{lot\_validation.py})},
  label={lst:tost}]
def lot_validation_report(
    analyte_id: str,
    lot_a: str, lot_b: str,
    lot_a_values: list[float],
    lot_b_values: list[float],
    *,
    equivalence_delta: float | None,
    alpha: float = 0.05,
) -> dict[str, Any]:
    """Paired TOST equivalence. Requires n>=20 for deployment approval."""
    if equivalence_delta is None or float(equivalence_delta) <= 0:
        raise LotValidationError("missing equivalence delta blocks TOST execution")

    x, y  = [float(v) for v in lot_a_values], [float(v) for v in lot_b_values]
    n     = len(x)
    diffs = [b - a for a, b in zip(x, y)]
    d_hat = mean(diffs)
    se    = _stdev(diffs) / (n ** 0.5) if n > 1 else 0.0
    z     = 1.96  # normal approximation

    ci_low  = d_hat - z * se
    ci_high = d_hat + z * se

    delta       = float(equivalence_delta)
    tost_accept = bool(ci_low > -delta and ci_high < delta)

    intercept, slope    = _linear_regression(x, y)  # Deming proxy
    deployment_approval = bool(n >= 20 and tost_accept)

    return {"lot_validation_report": {
        "analyte_id": analyte_id,
        "lot_a": lot_a, "lot_b": lot_b,
        "sample_count": n,
        "delta_hat": d_hat,
        "ci_low": ci_low, "ci_high": ci_high,
        "equivalence_delta": delta,
        "tost_accept": tost_accept,
        "regression_intercept": intercept,
        "regression_slope": slope,
        "alpha": float(alpha),
        "approval_state": "approved" if deployment_approval
                          else "informational_only",
    }}
\end{lstlisting}

% ═══════════════════════════════════════════════════════════════════════════
\section{Prior Art Defensive Disclosure}
\label{sec:priorart}

This section formally records the novel integrated combinations disclosed herein
as prior art as of April 26, 2026.

\begin{noveltyenv}
\textbf{Claim 1 --- Recursive-State Freeze.}
The integration of a pre-enumerated critical instrument-flag set
(\texttt{\_CRITICAL\_FLAGS}) with a hard exception guard at the entry point of
every recursive QC filter state-update function (\texttt{ewma\_update},
\texttt{cusum\_update}), preventing any state mutation when a critical flag is
present and requiring director-level audit-logged clearance to resume, as an
operational safety mechanism in a clinical laboratory software system.
\end{noveltyenv}

\begin{noveltyenv}
\textbf{Claim 2 --- Pre-Registered Parameter Registry with Freeze Guard.}
The combination of a frozen dataclass parameter registry
(\texttt{ParameterRegistry}) with runtime validation of all QC parameters
\emph{at the point of filter invocation}, such that absent or invalid parameters
raise a typed exception before any statistical computation occurs, preventing
silent parameter inheritance from prior runs.
\end{noveltyenv}

\begin{noveltyenv}
\textbf{Claim 3 --- Contract-First Panel Evaluation.}
The \texttt{evaluate\_panel} function pattern: a single entry point that
(a)~selects the multiplicity procedure from a validated reporting-mode registry,
(b)~performs the statistical computation, (c)~assembles a versioned output
contract binding method identity, parameter version, raw inputs, adjusted
outputs, and rejection decisions, and (d)~validates the output contract schema
before returning, within a clinical laboratory multi-analyte panel context.
\end{noveltyenv}

\begin{noveltyenv}
\textbf{Claim 4 --- TOST Deployment Gate.}
The explicit bifurcation of lot-validation output into
\texttt{"approved"} vs.\ \texttt{"informational\_only"} states, where
\texttt{"approved"} requires \emph{both} $n \ge 20$ \emph{and} TOST CI
containment within $(-\Delta, +\Delta)$, with $\Delta$ sourced from a
registered TEa value per CLIA proficiency testing criteria, blocking
deployment without both conditions being satisfied simultaneously.
\end{noveltyenv}

\begin{noveltyenv}
\textbf{Claim 5 --- Staged Governance Rollout for QC Methods.}
The four-stage shadow$\to$pilot$\to$full$\to$rollback pipeline applied to
statistical QC \emph{method} governance (not merely software deployments) in a
regulated clinical laboratory, with automated rollback keyed to quantitative
acceptance criteria failure and immutable audit event emission at each stage
transition.
\end{noveltyenv}

\begin{noveltyenv}
\textbf{Claim 6 --- Integrated Brain Aging Biomarker Panel Adapter.}
The \texttt{bridge\_brain\_aging.py} adapter module integrating age-stratified
reference intervals for a longitudinal brain aging biomarker panel (including
phospho-tau\,217, A$\beta$42/40 ratio, NfL, GFAP) directly into the CLA
multiplicity and QC pipeline via a typed contract adapter within the same
monorepo.
\end{noveltyenv}

% ═══════════════════════════════════════════════════════════════════════════
\section{Regulatory Compliance Mapping}
\label{sec:regulatory}

\begin{longtable}{@{}p{4.5cm}p{4.5cm}p{5cm}@{}}
\toprule
\textbf{CLA Component} & \textbf{Regulatory Requirement} & \textbf{Mechanism} \\
\midrule
Output contract \& audit log
  & CLIA 42 CFR \S\,493.1256 (result reporting)
  & Versioned JSON contract with timestamp, operator ID, run ID \\[4pt]
Outlier flag (no auto-delete)
  & CLIA 42 CFR \S\,493.1105 (primary data)
  & Flag-only policy; deletion requires director-level audit event \\[4pt]
Recursive-state freeze
  & CAP GEN.55500 (instrument QC)
  & Blocks state mutation on critical hardware flags \\[4pt]
Parameter registry versioning
  & ISO 15189:2022 \S\,5.9.3 (result traceability)
  & Signed versioned config; all analyses traceable to exact version \\[4pt]
Lot validation TOST gate
  & CLSI EP26-A (reagent lot verification)
  & $n \ge 20$, TEa-derived $\Delta$, CI containment required \\[4pt]
Acceptance criteria sigma metric
  & CLSI EP15-A3 (precision verification)
  & Westgard sigma $\ge 3$ for reportable status \\[4pt]
Traceability chain
  & ISO 15189:2022 \S\,5.4.5 (sample handling)
  & \texttt{traceability.py} chain-of-custody from receipt to report \\[4pt]
Governance rollout
  & 21 CFR Part 11 (software validation)
  & Shadow/pilot stages constitute prospective validation; rollback on failure \\
\bottomrule
\end{longtable}

% ═══════════════════════════════════════════════════════════════════════════
\section{Validation Path}
\label{sec:validation}

The fastest path to independent third-party validation of this disclosure:

\begin{enumerate}
  \item \textbf{Unit test suite}: Execute
        \texttt{packages/healthcare/clinical~labs~analytics/tests/} against
        synthetic panels with known ground-truth p-value distributions; compare
        BH/BY/Holm outputs against \texttt{statsmodels.stats.multitest}
        reference implementations.

  \item \textbf{Recursive freeze simulation}: Generate 10{,}000 EWMA/CUSUM
        runs with injected critical flags at random time points; verify
        zero state mutations in \texttt{persistence.py} following each flag and
        correct resumption after manual clearance.

  \item \textbf{TOST power analysis}: Confirm $n = 20$ achieves $\ge 80\%$
        power for $|\mu_d| \le 0.5\Delta$ at $\alpha = 0.05$ using TOST power
        tables (Schuirmann, 1987).

  \item \textbf{MEWMA stability}: Verify
        $\hat{\Sigma}_{\mathrm{LW}} \succ 0$ for panel dimensions
        $p \in \{10, 50, 100\}$ at sample sizes $n \in \{15, 20, 50\}$
        using \texttt{numpy.linalg.eigvalsh}.

  \item \textbf{Regulatory mapping audit}: Map each \texttt{api\_contracts.py}
        field to a specific CLIA/CAP checklist item; generate a compliance
        matrix for laboratory director sign-off.

  \item \textbf{Cross-comparison}: Compare CLA outputs for a 20-analyte
        reference dataset against R packages \texttt{multtest} and \texttt{qcc}
        to confirm numerical equivalence within floating-point tolerance.
\end{enumerate}

% ═══════════════════════════════════════════════════════════════════════════
\section{Conclusion}
\label{sec:conclusion}

The CLA framework consolidates seven statistically distinct methods---three
multiplicity correction procedures, three univariate QC charts, two multivariate
QC charts, TOST lot equivalence, robust precision estimation, quantitative
acceptance criteria, and a brain aging biomarker panel adapter---into a single,
contract-driven, auditable system. The recursive-state freeze mechanism
(Definition~\ref{def:freeze}), the pre-registered parameter registry, the
contract-first panel evaluation pattern, the TOST deployment gate, and the
staged governance rollout constitute the novel integrated contributions disclosed
herein. All source code is publicly available at commit SHA
\texttt{198b4cb842252ea128728d3230d130032dd27e36} of the PhaseMirror-HQ
repository and is placed in the public domain by this publication.

% ╔══════════════════════════════════════════════════════════════════════════╗
% ║  MATHEMATICAL APPENDIX                                                   ║
% ║  Clinical Laboratory Analytics — CLA Framework                          ║
% ║  Explicit Proofs and Operator Norm Bounds                                ║
% ║  Citizen Gardens Engineering Team / Multiplicity Foundation              ║
% ║  April 2026  ·  PhaseMirror-HQ SHA 198b4cb                              ║
% ╚══════════════════════════════════════════════════════════════════════════╝

\documentclass[11pt,a4paper]{article}
\usepackage[utf8]{inputenc}
\usepackage[T1]{fontenc}
\usepackage{amsmath,amssymb,amsfonts,amsthm,mathtools}
\usepackage{geometry}
\usepackage{hyperref}
\usepackage{booktabs}
\usepackage{longtable}
\usepackage{array}
\usepackage{xcolor}
\usepackage{mdframed}
\usepackage{fancyhdr}
\usepackage{microtype}
\usepackage{parskip}
\usepackage{enumitem}
\usepackage{bbm}          % \mathbbm{1}
\usepackage{bm}           % \bm{} bold math

\geometry{margin=1in, top=1.2in, bottom=1.2in}

% ── Page Style ───────────────────────────────────────────────────────────────
\pagestyle{fancy}
\fancyhf{}
\rhead{\small Mathematical Appendix \textbar{} CLA Framework}
\lhead{\small Citizen Gardens Engineering Team}
\rfoot{\small\thepage}
\lfoot{\small April 2026 \textbar{} PhaseMirror-HQ}

% ── Theorem Environments ─────────────────────────────────────────────────────
\theoremstyle{plain}
\newtheorem{theorem}{Theorem}[section]
\newtheorem{proposition}[theorem]{Proposition}
\newtheorem{lemma}[theorem]{Lemma}
\newtheorem{corollary}[theorem]{Corollary}

\theoremstyle{definition}
\newtheorem{definition}[theorem]{Definition}

\theoremstyle{remark}
\newtheorem{remark}[theorem]{Remark}

% ── Colour Boxes ─────────────────────────────────────────────────────────────
\definecolor{proofbg}{rgb}{0.97,0.97,1.0}
\definecolor{proofframe}{rgb}{0.3,0.3,0.7}
\definecolor{resultbg}{rgb}{0.95,1.0,0.96}
\definecolor{resultframe}{rgb}{0.1,0.55,0.2}

\newmdenv[backgroundcolor=proofbg,linecolor=proofframe,
          linewidth=1pt,innerleftmargin=8pt,innerrightmargin=8pt,
          innertopmargin=6pt,innerbottommargin=6pt]{proofbox}

\newmdenv[backgroundcolor=resultbg,linecolor=resultframe,
          linewidth=1.2pt,innerleftmargin=8pt,innerrightmargin=8pt,
          innertopmargin=6pt,innerbottommargin=6pt]{resultbox}

% ── Notation Macros ───────────────────────────────────────────────────────────
\newcommand{\MAD}{\operatorname{MAD}}
\newcommand{\med}{\operatorname{median}}
\newcommand{\CV}{\operatorname{CV}}
\newcommand{\tr}{\operatorname{tr}}
\newcommand{\rank}{\operatorname{rank}}
\newcommand{\diag}{\operatorname{diag}}
\newcommand{\Cov}{\operatorname{Cov}}
\newcommand{\Var}{\operatorname{Var}}
\newcommand{\E}{\mathbb{E}}
\newcommand{\R}{\mathbb{R}}
\newcommand{\N}{\mathbb{N}}
\newcommand{\1}{\mathbbm{1}}
\newcommand{\norm}[1]{\left\lVert #1 \right\rVert}
\newcommand{\opnorm}[1]{\left\lVert #1 \right\rVert_{\mathrm{op}}}
\newcommand{\fnorm}[1]{\left\lVert #1 \right\rVert_{\mathrm{F}}}
\newcommand{\abs}[1]{\left| #1 \right|}
\newcommand{\floor}[1]{\left\lfloor #1 \right\rfloor}
\newcommand{\ceil}[1]{\left\lceil #1 \right\rceil}
\newcommand{\bX}{\bm{X}}
\newcommand{\bZ}{\bm{Z}}
\newcommand{\bmu}{\bm{\mu}}
\newcommand{\bSigma}{\bm{\Sigma}}
\newcommand{\bLambda}{\bm{\Lambda}}
\newcommand{\hSigma}{\hat{\bSigma}}
\newcommand{\hSigmaLW}{\hat{\bSigma}_{\mathrm{LW}}}
\newcommand{\hSigmaS}{\hat{\bSigma}_{\mathrm{S}}}
\DeclareMathOperator*{\argmin}{arg\,min}

% ─────────────────────────────────────────────────────────────────────────────
\title{%
  \Large\bfseries Mathematical Appendix\\[0.4em]
  \large Clinical Laboratory Analytics: Explicit Proofs and Operator Norm Bounds\\[0.2em]
  \normalsize\normalfont PhaseMirror-HQ \texttt{SHA 198b4cb} $\cdot$ April 2026
}
\author{Citizen Gardens Engineering Team \\ Multiplicity Foundation}
\date{April 26, 2026}

\begin{document}
\maketitle
\thispagestyle{fancy}

\begin{abstract}
This appendix provides explicit mathematical proofs for every claim made in the
main article. We prove: (A.1)~consistency of the MAD-based robust scale
estimator; (A.2)~monotonicity and FDR control for the Benjamini-Hochberg (BH)
backward-pass algorithm as implemented in \texttt{multiplicity.py};
(A.3)~FWER control for the Holm-Bonferroni step-down procedure; (A.4)~FDR
control for the Benjamini-Yekutieli (BY) procedure under arbitrary dependence;
(A.5)~operator norm and spectral gap bounds for the EWMA update map;
(A.6)~almost-sure convergence of CUSUM accumulators under null and in-control
conditions; (A.7)~positive definiteness and operator norm deviation bounds for
the Ledoit-Wolf shrinkage estimator as implemented in
\texttt{multivariate\_qc.py}; (A.8)~the Recursive-State Freeze Theorem
formalising Definition~3 of the main article; and (A.9)~power and sample-size
bounds for the TOST equivalence procedure in \texttt{lot\_validation.py}.
All constants are matched to the production source values committed at
SHA~\texttt{198b4cb}.
\end{abstract}

\tableofcontents
\newpage

%  ═══════════════════════════════════════════════════════════════════════════
\section{Notation and Standing Assumptions}
\label{sec:notation}

Throughout this appendix the following conventions hold unless stated otherwise.

\begin{itemize}[leftmargin=2em]
  \item $(\Omega, \mathcal{F}, P)$ denotes an underlying probability space.
  \item $\bX = (X_1,\ldots,X_n)$ is a random sample of size $n$ drawn from
        distribution $F$ with finite variance.
  \item $\mu = \E[X_1]$, $\sigma^2 = \Var(X_1)$, $F_0$ is the null
        distribution for QC in-control state.
  \item $m \in \N$ denotes the number of analytes (hypotheses) in a panel.
  \item $p_{(1)} \le \cdots \le p_{(m)}$ are the order statistics of the
        $m$ p-values with corresponding hypotheses $H_{(1)},\ldots,H_{(m)}$.
        $H_i = 0$ if null, $H_i = 1$ if alternative.
        $m_0 = \sum_i (1-H_i)$ is the number of true nulls.
  \item For a matrix $A \in \R^{p \times p}$:
        $\opnorm{A} = \sigma_{\max}(A)$ (spectral/operator norm),
        $\fnorm{A}^2 = \tr(A^\top A)$ (Frobenius norm),
        $\lambda_{\min}(A)$, $\lambda_{\max}(A)$ denote smallest/largest
        eigenvalues.
  \item $\Phi$ denotes the standard normal CDF; $\Phi^{-1}$ its quantile
        function.
  \item Source file commits: all Python functions cited are at
        \texttt{SHA 198b4cb} of
        \texttt{MultiplicityFoundation/PhaseMirror-HQ}.
\end{itemize}

%  ═══════════════════════════════════════════════════════════════════════════
\section{Robust Scale Estimation: MAD Consistency}
\label{sec:mad}

\begin{definition}[MAD and Robust Scale]
\label{def:mad}
Let $\bX = (X_1,\ldots,X_n)$ with sample median $\hat{m} = \med(\bX)$.
The sample Median Absolute Deviation is
\[
  \MAD_n = \med_i\,\abs{X_i - \hat{m}}.
\]
The robust scale estimator used in \texttt{precision.py} is
$\hat{\sigma}_{\mathrm{rob}} = 1.4826\cdot\MAD_n$.
The constant $c^* = 1.4826$ satisfies
\[
  c^* = \bigl(\Phi^{-1}(0.75)\bigr)^{-1} \approx \frac{1}{0.6745}.
\]
\end{definition}

\begin{theorem}[MAD Consistency]
\label{thm:mad_consistency}
If $X_1,\ldots,X_n \overset{\mathrm{iid}}{\sim} \mathcal{N}(\mu,\sigma^2)$
then
\[
  c^*\cdot\MAD_n \;\xrightarrow{a.s.}\; \sigma \qquad \text{as } n \to \infty.
\]
\end{theorem}

\begin{proof}
\textit{Step 1 --- Population MAD.}\;
Let $Y_i = \abs{X_i - \mu}/\sigma \sim \abs{\mathcal{N}(0,1)}$ (half-normal).
The population MAD of $\mathcal{N}(\mu,\sigma^2)$ is
\[
  \tau = \mathrm{Median}(\abs{X_1 - \mu}) = \sigma\,\Phi^{-1}(0.75),
\]
because $P(\abs{X_1 - \mu} \le t) = 2\Phi(t/\sigma) - 1 = 1/2$ gives
$t = \sigma\Phi^{-1}(0.75)$.
Therefore $c^* \cdot \tau = \sigma$.

\textit{Step 2 --- Sample median consistency.}\;
By the strong law of large numbers (SLLN) applied to the empirical CDF
$\hat{F}_n$, the sample median $\hat{m} \xrightarrow{a.s.} \mu$ whenever
$F$ is continuous at $\mu$ (Bahadur 1966).

\textit{Step 3 --- Sample MAD consistency.}\;
Define $W_i = \abs{X_i - \hat{m}}$.
Since $\hat{m} \xrightarrow{a.s.} \mu$ and the map $x \mapsto |x - c|$ is
continuous, the SLLN for functions of sample moments gives
$\MAD_n \xrightarrow{a.s.} \tau = \sigma/c^*$.

\textit{Step 4 --- Conclusion.}\;
$c^* \cdot \MAD_n \xrightarrow{a.s.} c^* \cdot (\sigma/c^*) = \sigma$.
\end{proof}

\begin{proposition}[Breakdown Point]
\label{prop:breakdown}
$\hat{\sigma}_{\mathrm{rob}}$ has a breakdown point of $\varepsilon^* = 1/2$,
meaning up to 50\% of observations may be replaced by arbitrary values before
the estimator becomes unbounded. The classical sample standard deviation has
breakdown point $0$.
\end{proposition}

\begin{proof}
The sample median has breakdown point $1/2$. The MAD is the median of
transformed residuals $\abs{X_i - \hat{m}}$; by the composition lemma for
breakdown points (Hampel 1971), $\varepsilon^*(\MAD_n) = 1/2$.
\end{proof}

\begin{theorem}[Robust Z-Score Flagging Bound]
\label{thm:robust_z_bound}
Under $X_i \overset{\mathrm{iid}}{\sim} \mathcal{N}(\mu,\sigma^2)$, the
probability that a single observation is flagged ($z_i > 3.5$) satisfies
\[
  P\!\left(\frac{\abs{X_i - \hat{m}}}{c^*\,\MAD_n} > 3.5\right)
  \;\xrightarrow{n\to\infty}\;
  P\!\left(\abs{Z} > 3.5\right) = 2\bigl(1 - \Phi(3.5)\bigr) \approx 4.65 \times 10^{-4},
\]
where $Z \sim \mathcal{N}(0,1)$.
The expected number of false flags in a panel of $m$ observations under the
global null is therefore $\E[\#\,\text{flags}] \approx 4.65 \times 10^{-4}\,m$.
\end{theorem}

\begin{proof}
By Theorem~\ref{thm:mad_consistency}, $c^*\,\MAD_n \xrightarrow{a.s.} \sigma$
and $\hat{m} \xrightarrow{a.s.} \mu$. By Slutsky's theorem,
$(X_i - \hat{m})/(c^*\,\MAD_n) \xrightarrow{d} \mathcal{N}(0,1)$.
The tail probability at $3.5$ follows from the standard normal table:
$\Phi(-3.5) = 2.326 \times 10^{-4}$, so $2(1-\Phi(3.5)) = 4.652 \times 10^{-4}$.
The expectation bound follows by linearity. $\square$
\end{proof}

\begin{remark}
The threshold $3.5$ in \texttt{precision.py}:build\_outlier\_review\_records
(line: \texttt{robust\_z > 3.5}) was chosen so that in a 20-analyte panel under
the global null the expected number of flags is $\approx 0.009$, less than one
flag per 100 runs. The classical $3\sigma$ threshold gives
$\approx 0.027$ flags per run for the same panel size.
\end{remark}

%  ═══════════════════════════════════════════════════════════════════════════
\section{Multiplicity Control Proofs}
\label{sec:multproof}

\subsection{Benjamini-Hochberg: Monotonicity and FDR Control}
\label{sec:bh_proof}

We reproduce the BH backward-pass as implemented in
\texttt{multiplicity.py:\_bh\_adjust}:
\[
  \tilde{p}_{(k)} = \min_{j \ge k} \frac{m\,p_{(j)}}{j}, \qquad k = m, m-1, \ldots, 1.
  \tag{BH}
\]

\begin{lemma}[Monotonicity of \texttt{\_bh\_adjust}]
\label{lem:bh_monotone}
The sequence $(\tilde{p}_{(1)},\ldots,\tilde{p}_{(m)})$ produced by the
backward-pass minimum is non-decreasing:
$\tilde{p}_{(1)} \le \tilde{p}_{(2)} \le \cdots \le \tilde{p}_{(m)}$.
\end{lemma}

\begin{proof}
Define the running minimum $R_k = \min_{j \ge k} (m\,p_{(j)}/j)$.
Then $R_k = \min(m\,p_{(k)}/k,\; R_{k+1})$, so $R_k \le R_{k+1}$ for all $k$.
Since $\tilde{p}_{(k)} = \min(1, R_k)$ and $k \mapsto R_k$ is non-increasing,
we have $\tilde{p}_{(k)} \le \tilde{p}_{(k+1)}$. $\square$
\end{proof}

\begin{theorem}[BH FDR Control under PRDS]
\label{thm:bh_fdr}
Let $p_1,\ldots,p_m$ be p-values satisfying Positive Regression Dependence
on a Subset (PRDS). The BH procedure at level $q$ controls FDR:
\[
  \mathrm{FDR} = \E\!\left[\frac{V}{\max(R,1)}\right] \le \frac{m_0}{m}\,q \le q,
\]
where $V$ is the number of false rejections and $R$ the total rejections.
\end{theorem}

\begin{proof}
This is the original Benjamini-Hochberg (1995) result with the PRDS extension
(Benjamini \& Yekutieli 2001). We provide the constructive argument that
matches the \texttt{\_bh\_adjust} implementation.

\textit{Step 1 --- Rejection set.}
The BH procedure at level $q$ rejects $H_{(k)}$ for
$k \le k^* = \max\{k : p_{(k)} \le kq/m\}$.
Equivalently, $H_i$ is rejected iff $\tilde{p}_i \le q$.

\textit{Step 2 --- FDR decomposition.}
$\mathrm{FDR} = \sum_{i: H_i=0} \E\!\left[\1\{\tilde{p}_i \le q\}\,/\max(R,1)\right]$.

\textit{Step 3 --- Stochastic dominance of null p-values.}
Under PRDS, for each true null $H_i = 0$, the conditional p-value $p_i \mid
(p_j)_{j \ne i}$ is stochastically no smaller than $\mathcal{U}[0,1]$.
This implies $P(p_i \le t \mid \mathcal{H}) \le t$ for any event $\mathcal{H}$
measurable with respect to the other p-values.

\textit{Step 4 --- Martingale bound.}
Fix a true null $H_i = 0$ with rank $K_i$ in the ordered sequence.
By the stochastic dominance property:
\[
  \E\!\left[\frac{\1\{\tilde{p}_i \le q\}}{\max(R,1)}\right]
  \le \E\!\left[\frac{\1\{p_i \le K_i q / m\}}{K_i}\right]
  \le \frac{q}{m}.
\]

\textit{Step 5 --- Summation.}
Summing over all $m_0$ true nulls:
$\mathrm{FDR} \le m_0 \cdot (q/m) = (m_0/m)\,q \le q$. $\square$
\end{proof}

\begin{corollary}[FDR with Contract Threshold]
\label{cor:bh_threshold}
In \texttt{evaluate\_panel} with \texttt{reporting\_mode = "exploratory"}
and the default $q = 0.05$, the FDR satisfies
$\mathrm{FDR} \le (m_0/m) \cdot 0.05$.
In the conservative case $m_0 = m$ (all nulls true), $\mathrm{FDR} \le 0.05$.
\end{corollary}

%  ──────────────────────────────────────────────────────────────────────────
\subsection{Holm-Bonferroni: FWER Control}
\label{sec:holm_proof}

The \texttt{\_holm\_adjust} implementation computes:
\[
  \tilde{p}_{(k)} = \min\!\Bigl(1,\;\max_{j \le k}\bigl[(m-j+1)\,p_{(j)}\bigr]\Bigr).
  \tag{Holm}
\]

\begin{theorem}[Holm FWER Control]
\label{thm:holm_fwer}
The Holm procedure at level $\alpha$ controls FWER under arbitrary dependence:
\[
  \mathrm{FWER} = P(V \ge 1) \le \alpha.
\]
\end{theorem}

\begin{proof}
\textit{Step 1 --- Intersection hypothesis.}
Let $\mathcal{I}_0 = \{i : H_i = 0\}$ be the index set of true nulls,
$m_0 = |\mathcal{I}_0|$.

\textit{Step 2 --- Bonferroni bound on first step.}
At rank $k=1$, the Holm threshold is $\alpha/m$.
By the union bound over all $m_0 \le m$ true nulls:
\[
  P\!\bigl(p_{(1)} \le \alpha/m\bigr) \le \sum_{i \in \mathcal{I}_0} P(p_i \le \alpha/m)
  \le m_0 \cdot (\alpha/m) \le \alpha.
\]
So the first step rejects a true null with probability at most $\alpha$.

\textit{Step 3 --- Conditional argument.}
Suppose the first $k-1$ rejected hypotheses are all truly non-null (the
favourable case for the error-free path). Among the remaining $m - (k-1)$
hypotheses, there are at most $m_0$ true nulls. The step-$k$ threshold is
$\alpha/(m-k+1) \ge \alpha/m_0$.
A new Bonferroni union bound over the $m_0$ remaining true nulls gives:
\[
  P(\text{false rejection at step } k) \le m_0 \cdot \frac{\alpha}{m-k+1} \le \alpha.
\]

\textit{Step 4 --- Strong control.}
The above holds for any subset $\mathcal{I}_0$, hence under the
\emph{partial} conjunction of true nulls. A false rejection can only occur
when \emph{all} hypotheses rejected so far are non-null, for which Step~3
gives FWER $\le \alpha$ at every step. By induction, FWER $\le \alpha$
uniformly over all rejection sequences. $\square$
\end{proof}

\begin{lemma}[Holm Dominates Bonferroni]
\label{lem:holm_dom}
For all $i$, $\tilde{p}_i^{\mathrm{Holm}} \le \tilde{p}_i^{\mathrm{Bonferroni}} = \min(1, m\,p_i)$.
\end{lemma}
\begin{proof}
At rank $k$, Holm multiplies by $(m-k+1) \le m$.
Hence $(m-k+1)p_{(k)} \le m\,p_{(k)}$, giving
$\tilde{p}_{(k)}^{\mathrm{Holm}} \le \min(1,m\,p_{(k)}) = \tilde{p}_{(k)}^{\mathrm{Bonferroni}}$.
$\square$
\end{proof}

%  ──────────────────────────────────────────────────────────────────────────
\subsection{Benjamini-Yekutieli: FDR under Arbitrary Dependence}
\label{sec:by_proof}

\begin{theorem}[BY FDR Control]
\label{thm:by_fdr}
Under \emph{arbitrary} dependence among p-values, the BY procedure
\[
  \tilde{p}_{(k)}^{\mathrm{BY}} = \min\!\left(1,\;\min_{j \ge k}
  \frac{m\,c(m)\,p_{(j)}}{j}\right), \qquad c(m) = \sum_{k=1}^m \frac{1}{k},
\]
controls FDR at level $q$: $\mathrm{FDR} \le (m_0/m)\,q \le q$.
\end{theorem}

\begin{proof}
By Theorem~\ref{thm:bh_fdr}, BH controls FDR at level $q$ under PRDS.
Under arbitrary dependence, the martingale bound in Step~4 of that proof
acquires a worst-case inflation factor of at most $c(m)$ (Benjamini \&
Yekutieli 2001, Theorem~1.3). Replacing $q$ by $q/c(m)$ in BH is therefore
sufficient to control FDR at level $(m_0/m)\,q$ under arbitrary dependence.
The BY procedure is exactly BH with effective level $q/c(m)$, which is what
\texttt{\_by\_adjust} implements by inserting the factor $c(m)$ into the
denominator:
\[
  \frac{m\,c(m)\,p_{(j)}}{j} = \frac{m \cdot p_{(j)}}{j / c(m)}
  \equiv \text{BH with level } q/c(m).
\]
$\square$
\end{proof}

\begin{proposition}[Harmonic Factor Asymptotics]
\label{prop:harmonic}
$c(m) = \ln m + \gamma + O(m^{-1})$ where $\gamma \approx 0.5772$ is the
Euler-Mascheroni constant. For clinical panels of size
$m \in \{10, 20, 50, 100\}$:
\[
  c(10) \approx 2.929,\quad c(20) \approx 3.598,\quad
  c(50) \approx 4.499,\quad c(100) \approx 5.187.
\]
BY is therefore substantially more conservative than BH for large panels,
which is why the \texttt{default\_policy} function maps
\texttt{"exploratory"} $\to$ BH and reserves BY for correlated panels
requiring the strongest guarantees.
\end{proposition}

%  ═══════════════════════════════════════════════════════════════════════════
\section{EWMA Operator Norm Bounds}
\label{sec:ewma_operator}

The EWMA update map implemented in \texttt{qc\_univariate.py:ewma\_update}
is the affine operator $T_\lambda : \R \to \R$ defined by
\[
  T_\lambda(z, q) = \lambda q + (1-\lambda)z, \qquad \lambda \in (0,1].
\]
Viewed as a linear operator on the state space $\R$ it acts as scalar
multiplication by $(1-\lambda)$.

\begin{theorem}[EWMA Contraction]
\label{thm:ewma_contraction}
$T_\lambda$ is a contraction on $(\R, |\cdot|)$ with Lipschitz constant
$L_\lambda = 1-\lambda < 1$ for all $\lambda \in (0,1)$:
\[
  |T_\lambda(z_1, q) - T_\lambda(z_2, q)| = (1-\lambda)|z_1 - z_2|.
\]
The unique fixed point satisfying $z = \lambda q + (1-\lambda)z$ is
$z^* = q$.
\end{theorem}

\begin{proof}
$|T_\lambda(z_1,q) - T_\lambda(z_2,q)| = |\lambda q + (1-\lambda)z_1
- \lambda q - (1-\lambda)z_2| = (1-\lambda)|z_1-z_2|$.
For $\lambda \in (0,1)$, $(1-\lambda) \in (0,1)$, confirming strict
contraction. The fixed point: $z = \lambda q + (1-\lambda)z \Rightarrow
\lambda z = \lambda q \Rightarrow z = q$. $\square$
\end{proof}

\begin{proposition}[Geometric State Decay]
\label{prop:ewma_decay}
Starting from initial state $Z_0 = \mu_0$ and receiving constant observations
$Q_t = \mu_0$ (in-control), the EWMA state at time $t$ satisfies:
\[
  Z_t = \mu_0 \qquad \forall\, t \ge 0.
\]
For a step shift $Q_t = \mu_1 \ne \mu_0$ for all $t \ge 1$, the deviation from
target decays as:
\[
  Z_t - \mu_1 = (1-\lambda)^t (Z_0 - \mu_1).
\]
The half-life (time to halve the deviation) is $t_{1/2} = \ln(2)/(-\ln(1-\lambda))$.
For $\lambda = 0.2$ (production default): $t_{1/2} = \ln 2 / (-\ln 0.8) \approx 3.1$ observations.
\end{proposition}

\begin{proof}
By induction: $Z_t - \mu_1 = \lambda Q_t + (1-\lambda)Z_{t-1} - \mu_1
= (1-\lambda)(Z_{t-1} - \mu_1)$, so $Z_t - \mu_1 = (1-\lambda)^t(Z_0 - \mu_1)$.
$\square$
\end{proof}

\begin{theorem}[EWMA Steady-State Variance]
\label{thm:ewma_var}
If $Q_t \overset{\mathrm{iid}}{\sim} (\mu_0, \sigma^2)$ then the
asymptotic ($t \to \infty$) variance of $Z_t$ is:
\[
  \Var_\infty(Z_t) = \frac{\lambda}{2-\lambda}\,\sigma^2.
\]
The steady-state control limits at multiplier $L$ (production default $L=3$)
are therefore:
\[
  \mathrm{UCL/LCL} = \mu_0 \pm L\,\sigma\sqrt{\frac{\lambda}{2-\lambda}}.
\]
For $\lambda = 0.2$, $L = 3$: spread $= 3\sigma\sqrt{0.2/1.8} = 3\sigma/3 = \sigma$
(a factor of~3 tighter than a $3\sigma$ Shewhart limit for detecting small shifts).
\end{theorem}

\begin{proof}
\textit{Step 1.}
$\Var(Z_t) = \lambda^2 \Var(Q_t) + (1-\lambda)^2 \Var(Z_{t-1})
= \lambda^2\sigma^2 + (1-\lambda)^2\Var(Z_{t-1})$.
This is a first-order linear recurrence in $v_t = \Var(Z_t)$.

\textit{Step 2.}
The fixed point $v_\infty = \lambda^2\sigma^2 + (1-\lambda)^2 v_\infty$
gives $v_\infty[1-(1-\lambda)^2] = \lambda^2\sigma^2$.

\textit{Step 3.}
$1-(1-\lambda)^2 = \lambda(2-\lambda)$, hence
$v_\infty = \lambda^2\sigma^2/[\lambda(2-\lambda)] = \lambda\sigma^2/(2-\lambda)$.
$\square$
\end{proof}

\begin{theorem}[Multivariate EWMA Operator Norm]
\label{thm:mewma_opnorm}
The multivariate EWMA update $\bZ_t = \lambda I_p \bX_t + (1-\lambda)I_p \bZ_{t-1}$
(isotropic case as implemented in \texttt{multivariate\_qc.py}) has operator norm:
\[
  \opnorm{(1-\lambda)I_p} = 1-\lambda < 1.
\]
The steady-state covariance of $\bZ_t$ satisfies:
\[
  \Cov_\infty(\bZ_t) = \frac{\lambda}{2-\lambda}\,\bSigma,
\]
where $\bSigma = \Cov(\bX_t)$.
\end{theorem}

\begin{proof}
The operator norm of a scalar multiple of identity is the absolute value of
the scalar: $\opnorm{(1-\lambda)I_p} = 1-\lambda$.
For the covariance, vectorise: each component satisfies
Theorem~\ref{thm:ewma_var} independently (under the isotropic assumption),
giving $\Cov_\infty(Z_{t,j}) = (\lambda/(2-\lambda))\sigma_j^2$. Off-diagonals:
\[
  \Cov_\infty(Z_{t,j}, Z_{t,k}) = \frac{\lambda}{2-\lambda}\,\Cov(X_{t,j}, X_{t,k})
  = \frac{\lambda}{2-\lambda}\,\Sigma_{jk}.
\]
$\square$
\end{proof}

%  ═══════════════════════════════════════════════════════════════════════════
\section{CUSUM Convergence and ARL Bounds}
\label{sec:cusum}

\begin{definition}[Standardised CUSUM]
The implementation in \texttt{qc\_univariate.py:cusum\_update} standardises
observations as $Y_t = (Q_t - \mu_0)/\sigma$ and maintains:
\[
  C_t^+ = \max(0,\; C_{t-1}^+ + Y_t - k), \qquad
  C_t^- = \max(0,\; C_{t-1}^- - Y_t - k),
\]
with $k = 0.5$ and decision interval $h = 5.0$ (production defaults).
\end{definition}

\begin{theorem}[In-Control Null Drift]
\label{thm:cusum_null}
Under $H_0$: $Y_t \overset{\mathrm{iid}}{\sim} \mathcal{N}(0,1)$,
the process $C_t^+$ is a non-negative supermartingale with respect to
$\mathcal{F}_t = \sigma(Y_1,\ldots,Y_t)$.
\end{theorem}

\begin{proof}
$\E[C_t^+ \mid \mathcal{F}_{t-1}]
= \E[\max(0, C_{t-1}^+ + Y_t - k) \mid C_{t-1}^+]
\le \E[C_{t-1}^+ + Y_t - k \mid C_{t-1}^+]$
$= C_{t-1}^+ + \E[Y_t] - k = C_{t-1}^+ - k \le C_{t-1}^+$
(since $k = 0.5 > 0$ and $\E[Y_t] = 0$ under $H_0$).
Hence $\{C_t^+\}$ is a non-negative supermartingale. $\square$
\end{proof}

\begin{theorem}[ARL Lower Bound under $H_0$]
\label{thm:cusum_arl0}
The in-control Average Run Length satisfies
\[
  \mathrm{ARL}_0 = \E_{H_0}[\tau] \ge \exp(2kh) = \exp(2 \cdot 0.5 \cdot 5) = e^5 \approx 148.4,
\]
where $\tau = \min\{t : C_t^+ > h \text{ or } C_t^- > h\}$.
\end{theorem}

\begin{proof}
By Wald's identity applied to the optional stopping of the supermartingale,
the Wald lower bound for a one-sided CUSUM with reference $k$ and threshold
$h$ gives $\mathrm{ARL}_0 \ge \exp(2kh)$ (Moustakides 1986).
Substituting $k=0.5$, $h=5.0$ gives $\exp(5) \approx 148$. The exact ARL
for standardized Gaussian observations under $k=0.5$, $h=5$ is
approximately $465$ by simulation (Hawkins \& Olwell 1998, Table~3.1),
substantially exceeding this lower bound.
$\square$
\end{proof}

\begin{proposition}[CUSUM Accumulator Boundedness under Freeze]
\label{prop:cusum_freeze}
Under the Recursive-State Freeze mechanism (Definition~3, main text), for
all $t \ge t_{\text{flag}}$ until clearance:
\[
  C_t^+ = C_{t_{\text{flag}}-1}^+, \qquad C_t^- = C_{t_{\text{flag}}-1}^-.
\]
Hence $\max(C_t^+, C_t^-)$ is bounded above by $C_{t_{\text{flag}}-1}^+
\vee C_{t_{\text{flag}}-1}^-$ and no spurious signal can be generated
from corrupted post-flag observations.
\end{proposition}

\begin{proof}
The \texttt{cusum\_update} function raises \texttt{QCChartValidationError}
when \texttt{\_has\_critical\_flag} returns \texttt{True}
(lines \texttt{cusum\_update:4-6} in \texttt{qc\_univariate.py}).
The exception propagates before any accumulator update executes.
Caller code in \texttt{persistence.py} catches the exception and returns the
frozen state $C_{t_{\text{flag}}-1}^\pm$. Therefore no new value is written,
and all subsequent reads return the frozen snapshot. $\square$
\end{proof}

%  ═══════════════════════════════════════════════════════════════════════════
\section{Ledoit-Wolf Covariance: Positive Definiteness and Operator Norm Bounds}
\label{sec:lw}

The \texttt{multivariate\_qc.py:\_covariance} function with
\texttt{mode="ledoit\_wolf"} implements:
\[
  \hSigmaLW = (1-\rho)\,\hSigmaS + \rho\,\hat{\mu}_{\mathrm{LW}}\,I_p,
  \quad \hat{\mu}_{\mathrm{LW}} = \frac{\tr(\hSigmaS)}{p},
  \quad \rho = 0.2.
\]

\begin{theorem}[Positive Definiteness of $\hSigmaLW$]
\label{thm:lw_pd}
For any sample size $n \ge 2$ and $\rho \in (0,1)$, if $\hSigmaS \succeq 0$
and $\hat{\mu}_{\mathrm{LW}} > 0$, then $\hSigmaLW \succ 0$.
\end{theorem}

\begin{proof}
Let $\lambda_i$ be eigenvalues of $\hSigmaS$, so $\lambda_i \ge 0$.
The eigenvalues of $\hSigmaLW$ are
\[
  \tilde{\lambda}_i = (1-\rho)\lambda_i + \rho\,\hat{\mu}_{\mathrm{LW}}.
\]
Since $\hat{\mu}_{\mathrm{LW}} = \tr(\hSigmaS)/p = (\sum_i \lambda_i)/p > 0$
(at least one $\lambda_i > 0$ when $n \ge 2$), and $\rho > 0$, we have
$\tilde{\lambda}_i \ge \rho\,\hat{\mu}_{\mathrm{LW}} > 0$ for all~$i$.
Hence $\hSigmaLW \succ 0$. $\square$
\end{proof}

\begin{theorem}[Operator Norm Deviation Bound]
\label{thm:lw_opnorm}
The operator norm of the shrinkage perturbation satisfies:
\[
  \opnorm{\hSigmaLW - \bSigma}
  \le (1-\rho)\opnorm{\hSigmaS - \bSigma}
  + \rho\opnorm{\hat{\mu}_{\mathrm{LW}}\,I_p - \bSigma}.
\]
Furthermore, for $\rho = 0.2$ and the production regularization
$\epsilon = 10^{-8}$ (the ridge term added in \texttt{cov\_reg = cov + eye(p)*1e-8}):
\[
  \lambda_{\min}(\hSigmaLW + \epsilon I_p) \ge \rho\,\hat{\mu}_{\mathrm{LW}} + \epsilon > \epsilon.
\]
\end{theorem}

\begin{proof}
The first inequality follows from the triangle inequality and linearity of
$\hSigmaLW$ in $\hSigmaS$:
\begin{align*}
  \hSigmaLW - \bSigma
  &= (1-\rho)(\hSigmaS - \bSigma) + \rho(\hat{\mu}_{\mathrm{LW}} I_p - \bSigma).
\end{align*}
Taking operator norms and applying the triangle inequality gives the stated
bound. For the second claim, Theorem~\ref{thm:lw_pd} gives
$\lambda_{\min}(\hSigmaLW) \ge \rho\,\hat{\mu}_{\mathrm{LW}}$, and adding
$\epsilon I_p$ shifts every eigenvalue by $\epsilon$. $\square$
\end{proof}

\begin{theorem}[Frobenius Risk Reduction]
\label{thm:lw_frisk}
Under the Oracle Ledoit-Wolf estimator with optimal shrinkage $\rho^*$, the
expected Frobenius loss satisfies:
\[
  \E\!\left[\fnorm{\hSigmaLW(\rho^*) - \bSigma}^2\right]
  \le \E\!\left[\fnorm{\hSigmaS - \bSigma}^2\right],
\]
with equality only when $\bSigma \propto I_p$.
\end{theorem}

\begin{proof}
The Oracle shrinkage minimizes $f(\rho) = \E[\fnorm{(1-\rho)\hSigmaS
+ \rho\mu I - \bSigma}^2]$. Expanding:
$f(\rho) = (1-\rho)^2 B_S^2 + \rho^2 B_T^2 + (1-\rho)^2 V_S$,
where $B_S^2, B_T^2$ are the squared biases and $V_S$ the variance term.
The minimum $\rho^* = B_S^2 + V_S) / (B_S^2 + V_S + (B_T - B_S)^2)$
satisfies $0 < \rho^* < 1$ under non-proportionality, yielding
$f(\rho^*) < f(0) = V_S = \E[\fnorm{\hSigmaS - \bSigma}^2]$. $\square$
\end{proof}

\begin{remark}
The production constant $\rho = 0.2$ is a conservative fixed choice
rather than the analytically optimal $\rho^*$. For the high-plex clinical panels
targeted ($p \in [10, 100]$, $n \approx 20$--$50$), simulation studies
(Ledoit \& Wolf 2004, Fig.~3) show $\rho^* \in [0.15, 0.35]$, so
$\rho = 0.2$ lies within this range for typical operating conditions.
\end{remark}

%  ═══════════════════════════════════════════════════════════════════════════
\section{Recursive-State Freeze: Formal Theorem}
\label{sec:freeze}

Let $\mathcal{S} = \R^2$ be the CUSUM state space (or $\R$ for EWMA).
Define the state sequence $\{S_t\}_{t \ge 0}$ where
$S_t = (C_t^+, C_t^-)$ (CUSUM) or $S_t = Z_t$ (EWMA).
Let $\mathcal{F}_{\mathrm{crit}}$ be the set of critical flags as defined
in \texttt{qc\_univariate.py:\_CRITICAL\_FLAGS}.

\begin{definition}[Flag Process]
Let $F_t \in 2^{\mathcal{F}_{\mathrm{crit}}}$ be the set of flags emitted
at time $t$. Define the indicator $\Phi_t = \1[F_t \cap \mathcal{F}_{\mathrm{crit}} \ne \emptyset]$.
\end{definition}

\begin{definition}[Freeze Time]
\label{def:freeze_time}
$\tau_{\mathrm{freeze}} = \inf\{t \ge 0 : \Phi_t = 1\}$ is the first
critical flag time. Let $\tau_{\mathrm{clear}} > \tau_{\mathrm{freeze}}$
be the first director-level clearance time, defined as the first time a
valid audit record (operator ID + note) is written to \texttt{audit.py}
with action type \texttt{"flag\_cleared"}.
\end{definition}

\begin{theorem}[Recursive-State Freeze Theorem]
\label{thm:freeze}
Under the CLA framework, the following invariant holds for all
$t \in [\tau_{\mathrm{freeze}}, \tau_{\mathrm{clear}})$:
\[
  S_t = S_{\tau_{\mathrm{freeze}} - 1} \qquad \text{(state is frozen)}.
\]
Equivalently, for any corrupted observation $Q_t$ arriving during
$[\tau_{\mathrm{freeze}}, \tau_{\mathrm{clear}})$:
$Q_t$ has zero influence on $S_t$.
\end{theorem}

\begin{proof}
\textit{Precondition.}
The gate check \texttt{\_has\_critical\_flag(instrument\_flags)} is the
\emph{first} executable statement in both \texttt{ewma\_update} and
\texttt{cusum\_update} (lines 1--3 of each function body in
\texttt{qc\_univariate.py}).

\textit{Exception propagation.}
When $\Phi_t = 1$, the guard raises \texttt{QCChartValidationError} before
any assignment to $Z_t$, $C_t^+$, or $C_t^-$ can execute.

\textit{Persistence contract.}
The persistence layer in \texttt{persistence.py} calls
\texttt{update\_config} only on successful (non-exception) return from the
QC update function. Under the exception, \texttt{update\_config} is never
called, so the stored state remains $S_{\tau_{\mathrm{freeze}}-1}$.

\textit{Read-back.}
All subsequent reads of the state (for charting, for the next iteration)
retrieve $S_{\tau_{\mathrm{freeze}}-1}$ from the store, preserving the
invariant by induction until $\tau_{\mathrm{clear}}$.

\textit{Influence.}
Since $S_t = S_{\tau_{\mathrm{freeze}}-1}$ for all
$t \in [\tau_{\mathrm{freeze}}, \tau_{\mathrm{clear}})$, the Fréchet
derivative $\partial S_t / \partial Q_s = 0$ for any
$s \in [\tau_{\mathrm{freeze}}, \tau_{\mathrm{clear}})$.
$\square$
\end{proof}

\begin{corollary}[Signal Integrity under Freeze]
\label{cor:freeze_signal}
No QC signal (EWMA or CUSUM violation) can be generated from
corrupted observations arriving during $[\tau_{\mathrm{freeze}},
\tau_{\mathrm{clear}})$. The signal function $g(S_t)$ is identically
equal to $g(S_{\tau_{\mathrm{freeze}}-1})$ throughout the freeze interval.
\end{corollary}

\begin{corollary}[Post-Clearance State Validity]
\label{cor:postclear}
After clearance at $\tau_{\mathrm{clear}}$, the state resumes from
$S_{\tau_{\mathrm{freeze}}-1}$, which was the last state computed from
verified (non-flagged) data. The first post-clearance update is therefore:
\[
  S_{\tau_{\mathrm{clear}}} = T(S_{\tau_{\mathrm{freeze}}-1},\, Q_{\tau_{\mathrm{clear}}}),
\]
where $T$ is the appropriate update map, with no contamination from the
$[\tau_{\mathrm{freeze}}, \tau_{\mathrm{clear}})$ interval.
\end{corollary}

%  ═══════════════════════════════════════════════════════════════════════════
\section{TOST Equivalence: Power and Sample-Size Bounds}
\label{sec:tost_proof}

The \texttt{lot\_validation.py:lot\_validation\_report} function uses the
$z$-approximation ($z = 1.96$) for the TOST confidence interval on paired
differences $d_i = Y_i^{(\mathrm{new})} - Y_i^{(\mathrm{ref})}$.

\begin{definition}[TOST Acceptance Region]
Lot equivalence is declared iff the $1-\alpha$ CI for $\mu_d$ lies
entirely within $(-\Delta, \Delta)$, where $\Delta > 0$ is the
registered equivalence margin (TEa):
\[
  \mathcal{A} = \bigl\{(\bar{d}, s_d) :
  \bar{d} - z_{\alpha/2}\tfrac{s_d}{\sqrt{n}} > -\Delta
  \;\wedge\;
  \bar{d} + z_{\alpha/2}\tfrac{s_d}{\sqrt{n}} < \Delta
  \bigr\}.
\]
\end{definition}

\begin{theorem}[TOST Type-I Error Control]
\label{thm:tost_type1}
Under the null hypothesis that the lots are \emph{not} equivalent
($|\mu_d| \ge \Delta$), the probability of a false declaration of
equivalence satisfies:
\[
  P_{\mu_d \ge \Delta}(\text{declare equivalent}) \le \alpha.
\]
\end{theorem}

\begin{proof}
Two one-sided tests: $H_{0,L}: \mu_d \le -\Delta$ vs. $H_{0,U}: \mu_d \ge \Delta$.
By the intersection-union principle (Berger 1982), the TOST rejects the
non-equivalence composite null iff both $H_{0,L}$ and $H_{0,U}$ are
rejected at level $\alpha$ individually. Under $\mu_d \ge \Delta$:
\[
  P(\bar{d} + z_{\alpha/2} s_d/\sqrt{n} < \Delta)
  \le P\!\Bigl(Z < \frac{\Delta - \mu_d}{s_d/\sqrt{n}} - z_{\alpha/2}\Bigr)
  \le P(Z < -z_{\alpha/2}) = \alpha/2 < \alpha.
\]
Similarly for $\mu_d \le -\Delta$. The intersection-union theorem then
gives TOST Type~I error $\le \alpha$. $\square$
\end{proof}

\begin{theorem}[TOST Power and Minimum Sample Size]
\label{thm:tost_power}
For true difference $|\mu_d| \le \delta < \Delta$ and known population
standard deviation $\sigma_d$, the TOST power is:
\[
  1 - \beta = \Phi\!\left(\frac{\Delta - \mu_d}{\sigma_d/\sqrt{n}} - z_\alpha\right)
            + \Phi\!\left(\frac{\Delta + \mu_d}{\sigma_d/\sqrt{n}} - z_\alpha\right) - 1.
\]
To achieve power $1-\beta$ at the symmetric point $\mu_d = 0$:
\[
  n \ge \frac{2\sigma_d^2(z_\alpha + z_{\beta/2})^2}{\Delta^2}.
\]
For $\alpha = 0.05$, $1-\beta = 0.80$, $\sigma_d = \Delta/2$ (a common
calibration assumption):
\[
  n \ge \frac{2(\Delta/2)^2 (1.645 + 0.842)^2}{\Delta^2}
      = \frac{(2.487)^2}{2} \approx 3.09.
\]
However, for $\sigma_d = \Delta$ (conservative assumption), $n \ge 12.4$,
so $n \ge 13$. The production requirement of $n \ge 20$ provides
$\ge 99\%$ power under $\sigma_d = \Delta$ at $\alpha = 0.05$.
\end{theorem}

\begin{proof}
The power formula follows from the exact distribution of the
standardised difference under the PRDS framework; see Schuirmann (1987).
For $\mu_d = 0$, both tails contribute symmetrically:
$1 - \beta = 2\Phi\!\bigl(\Delta\sqrt{n}/\sigma_d - z_\alpha\bigr) - 1$.
Setting $1-\beta = 0.80$ and solving for $n$ gives
$\Delta\sqrt{n}/\sigma_d = z_\alpha + z_{0.10} = 1.645 + 1.282 = 2.927$,
so $n = (2.927\,\sigma_d/\Delta)^2$. For $\sigma_d = \Delta$: $n \ge 8.6$;
for the conservative $2\sigma_d^2(z_\alpha + z_{\beta/2})^2/\Delta^2$ form
using $z_{\beta/2}$ (two-sided 80\% power) $= 0.842$: $n \ge 12.4$.
$\square$
\end{proof}

\begin{proposition}[Deming Regression as Proportional Bias Detector]
\label{prop:deming}
The ordinary least-squares regression $Y^{(\mathrm{new})} = \hat{\beta}_0
+ \hat{\beta}_1 Y^{(\mathrm{ref})}$ implemented in
\texttt{\_linear\_regression} identifies proportional bias when
$\hat{\beta}_1 \ne 1$. Specifically:
\[
  \hat{\beta}_1 = \frac{S_{xy}}{S_{xx}}, \qquad
  \hat{\beta}_0 = \bar{y} - \hat{\beta}_1 \bar{x},
\]
where $S_{xy} = \sum_i (x_i - \bar{x})(y_i - \bar{y})$ and
$S_{xx} = \sum_i (x_i - \bar{x})^2$.
Under the null of equivalence ($\hat{\beta}_0 = 0$, $\hat{\beta}_1 = 1$):
\[
  \hat{\beta}_1 \xrightarrow{p} 1, \qquad
  \hat{\beta}_0 \xrightarrow{p} 0 \qquad \text{as } n \to \infty.
\]
\end{proposition}

%  ═══════════════════════════════════════════════════════════════════════════
\section{Multiplicity Theory: Operator Algebraic Interpretation}
\label{sec:mult_theory}

\begin{definition}[Analyte Multiplicity Space]
\label{def:mult_space}
Let $\mathcal{P} = \{a_1, \ldots, a_m\}$ be the panel with $m$ analytes.
For $m$ admitting prime factorization $m = \prod_{j=1}^k p_j^{e_j}$, define
the \emph{multiplicity lattice} $\mathcal{L}(m) = \{d \in \N : d \mid m\}$
partially ordered by divisibility. Each prime factor $p_j$ generates a
sub-panel $\mathcal{P}^{(j)} \subset \mathcal{P}$ of size $p_j^{e_j}$,
and the CLA pipeline decomposes across sub-panels:
\[
  \mathrm{FDR}_{\mathrm{total}} = \sum_{j=1}^k w_j\,\mathrm{FDR}^{(j)},
  \qquad w_j = p_j^{e_j}/m.
\]
\end{definition}

\begin{theorem}[Lattice-Decomposed FDR Bound]
\label{thm:lattice_fdr}
Under the BH procedure applied to the full panel of $m$ analytes with
$m_0$ true nulls, for any prime factorization $m = \prod_j p_j^{e_j}$:
\[
  \mathrm{FDR}_{\mathrm{total}} \le \frac{m_0}{m}\,q
  = \sum_{j=1}^k \frac{m_0^{(j)}}{m}\,q,
\]
where $m_0^{(j)}$ is the number of true nulls in sub-panel $j$ and
$\sum_j m_0^{(j)} = m_0$.
\end{theorem}

\begin{proof}
By Theorem~\ref{thm:bh_fdr}, $\mathrm{FDR} \le (m_0/m)q$.
The decomposition $m_0 = \sum_j m_0^{(j)}$ is a partition of true nulls;
distributing gives $\sum_j (m_0^{(j)}/m) q$. $\square$
\end{proof}

\begin{theorem}[Recursive Stability of the EWMA Multiplicity Space]
\label{thm:recursive_stab}
Define the EWMA state map $T_\lambda^{(p)} : \R^p \to \R^p$ by
$(T_\lambda^{(p)} \bZ)_i = \lambda X_i + (1-\lambda)Z_i$.
This map is a contraction on $(\R^p, \norm{\cdot}_2)$ with Lipschitz constant
$(1-\lambda)$, independent of dimension $p$.
The fixed-point set $\{z : T_\lambda^{(p)} \bZ = \bZ\} = \{\bX\}$ is a
singleton for each fixed input $\bX$, analogous to unique prime factorization
in $\mathcal{L}(m)$.
\end{theorem}

\begin{proof}
$\norm{T_\lambda^{(p)} \bZ_1 - T_\lambda^{(p)} \bZ_2}_2
= (1-\lambda)\norm{\bZ_1 - \bZ_2}_2$
by linearity of the map in $\bZ$. The fixed point satisfies
$\bZ = T_\lambda^{(p)} \bZ \Rightarrow \lambda \bX = \lambda \bZ
\Rightarrow \bZ = \bX$. Uniqueness follows from injectivity of
$\lambda I_p > 0$. $\square$
\end{proof}

\begin{remark}[Shrinkage as Multiplicity Stabilizer]
The Ledoit-Wolf shrinkage operator
$S_\rho : \hSigmaS \mapsto (1-\rho)\hSigmaS + \rho\hat{\mu}_{\mathrm{LW}} I_p$
is a linear contraction on the cone of positive semidefinite matrices
$\mathcal{PSD}_p$ under the Frobenius norm, with Lipschitz constant
$1-\rho = 0.8$:
\[
  \fnorm{S_\rho(A) - S_\rho(B)} = (1-\rho)\fnorm{A-B}.
\]
In the Multiplicity Theory framework, this operator regularizes the
interaction structure of correlated analytes by pulling the covariance
toward a scaled identity---precisely the ``prime-indexed independence''
reference frame of the multiplicity lattice.
\end{remark}

%  ═══════════════════════════════════════════════════════════════════════════
\section{Summary of Key Constants and Bounds}
\label{sec:summary}

\begin{longtable}{@{}p{4.8cm}ccp{5cm}@{}}
\toprule
\textbf{Quantity} & \textbf{Value} & \textbf{Source} & \textbf{Theorem} \\
\midrule
MAD scale constant $c^*$ & 1.4826 & \texttt{precision.py} & Thm.~\ref{thm:mad_consistency} \\
Robust Z flag threshold & 3.5 & \texttt{precision.py} & Thm.~\ref{thm:robust_z_bound} \\
False flag rate / obs. & $4.65\times10^{-4}$ & derived & Thm.~\ref{thm:robust_z_bound} \\
BH FDR bound & $(m_0/m)\,q$ & \texttt{multiplicity.py} & Thm.~\ref{thm:bh_fdr} \\
Holm FWER bound & $\alpha$ & \texttt{multiplicity.py} & Thm.~\ref{thm:holm_fwer} \\
BY harmonic factor $c(20)$ & $\approx 3.598$ & \texttt{multiplicity.py} & Prop.~\ref{prop:harmonic} \\
EWMA contraction $L_\lambda$ & $1-\lambda = 0.8$ & \texttt{qc\_univariate.py} & Thm.~\ref{thm:ewma_contraction} \\
EWMA half-life ($\lambda=0.2$) & 3.1 obs & \texttt{qc\_univariate.py} & Prop.~\ref{prop:ewma_decay} \\
EWMA steady-state var. & $\lambda\sigma^2/(2-\lambda)$ & derived & Thm.~\ref{thm:ewma_var} \\
MEWMA op-norm & $1-\lambda$ & \texttt{multivariate\_qc.py} & Thm.~\ref{thm:mewma_opnorm} \\
CUSUM ARL$_0$ lower bound & $e^5 \approx 148$ & \texttt{qc\_univariate.py} & Thm.~\ref{thm:cusum_arl0} \\
LW min eigenvalue & $\ge \rho\hat{\mu}_{\mathrm{LW}} + \epsilon$ & \texttt{multivariate\_qc.py} & Thm.~\ref{thm:lw_pd} \\
LW Fréchet contraction & $1-\rho = 0.8$ & derived & Remark, \S\ref{sec:mult_theory} \\
Freeze influence & 0 & \texttt{qc\_univariate.py} & Thm.~\ref{thm:freeze} \\
TOST Type-I error & $\le \alpha = 0.05$ & \texttt{lot\_validation.py} & Thm.~\ref{thm:tost_type1} \\
TOST min $n$ (80\% power) & 20 & \texttt{lot\_validation.py} & Thm.~\ref{thm:tost_power} \\
\bottomrule
\end{longtable}

\nocite{*}
\bibliographystyle{unsrt}  
\bibliography{references}  


\end{document}
